Large language models that pass standard safety benchmarks may still harbor context-dependent vulnerabilities that direct-prompt evaluations fail to detect.
We investigate whether adversarial chain-of-thought (CoT) prompting --- prompts designed to exploit multi-step reasoning contexts --- can surface hidden misalignment in models deemed robust under conventional evaluations.
We test five prompting conditions (direct, simple CoT, snowball CoT, context manipulation CoT, and decomposition CoT) across 80 harmful behaviors from AdvBench on two widely deployed models, GPT-4o-mini and GPT-4o.
Our central finding is that context manipulation CoT prompting, which frames harmful requests within an academic analytical context, increases the attack success rate (ASR) on GPT-4o-mini from 2.5\% to 20.0\% --- an 8$\times$ increase that is statistically significant after Bonferroni correction (McNemar's $p=0.0001$, Cohen's $h=0.61$).
GPT-4o shows the same directional pattern (0.0\% to 7.5\%, $p=0.031$) with a lower absolute effect.
Other CoT strategies, including snowball and decomposition prompting, do not significantly outperform the direct baseline.
Neither model exhibits over-refusal on safe XSTest prompts (0.0\%).
These results demonstrate that standard direct-prompt safety evaluations systematically underestimate model vulnerability and that safety assessments should incorporate adversarial reasoning probes, particularly context manipulation strategies, to provide more faithful estimates of deployment risk.
